% Options for packages loaded elsewhere
\PassOptionsToPackage{unicode}{hyperref}
\PassOptionsToPackage{hyphens}{url}
\PassOptionsToPackage{dvipsnames,svgnames,x11names}{xcolor}
%
\documentclass[
  10pt,
  a4paper,
]{article}
\usepackage{amsmath,amssymb}
\usepackage{setspace}
\usepackage{iftex}
\ifPDFTeX
  \usepackage[T1]{fontenc}
  \usepackage[utf8]{inputenc}
  \usepackage{textcomp} % provide euro and other symbols
\else % if luatex or xetex
  \usepackage{unicode-math} % this also loads fontspec
  \defaultfontfeatures{Scale=MatchLowercase}
  \defaultfontfeatures[\rmfamily]{Ligatures=TeX,Scale=1}
\fi
\usepackage{lmodern}
\ifPDFTeX\else
  % xetex/luatex font selection
  \setmainfont[]{Liberation Serif}
  \setsansfont[]{Liberation Sans}
  \setmonofont[]{DejaVu Sans Mono}
\fi
% Use upquote if available, for straight quotes in verbatim environments
\IfFileExists{upquote.sty}{\usepackage{upquote}}{}
\IfFileExists{microtype.sty}{% use microtype if available
  \usepackage[]{microtype}
  \UseMicrotypeSet[protrusion]{basicmath} % disable protrusion for tt fonts
}{}
\makeatletter
\@ifundefined{KOMAClassName}{% if non-KOMA class
  \IfFileExists{parskip.sty}{%
    \usepackage{parskip}
  }{% else
    \setlength{\parindent}{0pt}
    \setlength{\parskip}{6pt plus 2pt minus 1pt}}
}{% if KOMA class
  \KOMAoptions{parskip=half}}
\makeatother
\usepackage{xcolor}
\usepackage[margin=22mm]{geometry}
\setlength{\emergencystretch}{3em} % prevent overfull lines
\providecommand{\tightlist}{%
  \setlength{\itemsep}{0pt}\setlength{\parskip}{0pt}}
\setcounter{secnumdepth}{-\maxdimen} % remove section numbering
\ifLuaTeX
\usepackage[bidi=basic]{babel}
\else
\usepackage[bidi=default]{babel}
\fi
\babelprovide[main,import]{korean}
\ifPDFTeX
\else
\babelfont{rm}[]{Liberation Serif}
\fi
% get rid of language-specific shorthands (see #6817):
\let\LanguageShortHands\languageshorthands
\def\languageshorthands#1{}
\usepackage{xeCJK}
\setCJKmainfont{Noto Serif CJK KR}
\setCJKsansfont{Noto Sans CJK KR}
\setCJKmonofont{Noto Sans Mono CJK KR}
\xeCJKsetup{CJKspace=true}
\usepackage{xurl}
\usepackage{fvextra}
\usepackage{fancyhdr}
\AtBeginDocument{\DefineVerbatimEnvironment{verbatim}{Verbatim}{breaklines=true,breakanywhere=true,fontsize=\small}}
\setlength{\emergencystretch}{4em}
\setlength{\headheight}{15pt}
\pagestyle{fancy}\fancyhf{}
\fancyhead[L]{\small 중성자별 구조와 비방사 진동}
\fancyfoot[C]{\thepage}
\renewcommand{\headrulewidth}{0.3pt}
\setcounter{tocdepth}{2}
\ifLuaTeX
  \usepackage{selnolig}  % disable illegal ligatures
\fi
\usepackage{bookmark}
\IfFileExists{xurl.sty}{\usepackage{xurl}}{} % add URL line breaks if available
\urlstyle{same}
\hypersetup{
  pdftitle={A02 · 이분법과 실패를 보존하는 근 찾기},
  pdflang={ko-KR},
  colorlinks=true,
  linkcolor={Maroon},
  filecolor={Maroon},
  citecolor={Blue},
  urlcolor={Blue},
  pdfcreator={LaTeX via pandoc}}

\title{A02 · 이분법과 실패를 보존하는 근 찾기}
\usepackage{etoolbox}
\makeatletter
\providecommand{\subtitle}[1]{% add subtitle to \maketitle
  \apptocmd{\@title}{\par {\large #1 \par}}{}{}
}
\makeatother
\subtitle{과제 큐시트}
\author{}
\date{}

\begin{document}
\maketitle

\setstretch{1.12}
\textbf{선행:} A01　 \textbf{독립 작업:} 3--5시간 (계획치)

\subsection{작은 잔차는 근에 가까워졌다는
뜻일까?}\label{uxc791uxc740-uxc794uxcc28uxb294-uxadfcuxc5d0-uxac00uxae4cuxc6ccuxc84cuxb2e4uxb294-uxb73buxc77cuxae4c}

같은 함수 전체를 작은 양의 상수로 곱한다. 근의 위치, 잔차의 크기,
이분법의 구간 폭 중 무엇이 변하는지 계산 전에 나누어 적자.

먼저 한두 문장으로 답한 뒤 아래 설명을 읽자. 첫 예측이 틀렸거나 아직
모르겠다고 적는 것은 감점하지 않는다.

\subsection{필요한 관계와 읽을
위치}\label{uxd544uxc694uxd55c-uxad00uxacc4uxc640-uxc77duxc744-uxc704uxce58}

\(F(a)F(b)<0\)는 연속성 가정 아래에서 근 존재의 충분조건이다. 실제 부호
비교는 곱셈 없이 한다. \(b_k-a_k=(b_0-a_0)2^{-k}\).

\begin{itemize}
\tightlist
\item
  \path|notes/chapters/ch02_numerics.md|
\item
  \path|learner/src/numerics/ns_bisection.f90|
\end{itemize}

\subsection{계산으로
옮기기}\label{uxacc4uxc0b0uxc73cuxb85c-uxc62euxae30uxae30}

기본 경로는 제공 함수의 호출과 해석이다. 아래 \texttt{learner} 시험은
같은 물리 함수를 직접 작성한 경우에 실행한다. 제공 예제를 검사할 때에는
명시적으로 \texttt{calculation}을 선택한다.

\begin{enumerate}
\def\labelenumi{\arabic{enumi}.}
\tightlist
\item
  풀이에서 bracket의 연속성 전제와 갱신할 절반을 확인한다.
\item
  제공 근 찾기를 호출하고 정상 근·끝점 근·실패 상태를 비교한다.
\item
  함수의 배율을 바꿔 잔차 조건과 좌표 조건을 따로 읽는다. 직접 구현은
  같은 시험을 사용하는 선택 경로다.
\end{enumerate}

\newpage

\subsection{실행해서 무엇을 볼
것인가}\label{uxc2e4uxd589uxd574uxc11c-uxbb34uxc5c7uxc744-uxbcfc-uxac83uxc778uxac00}

패키지 루트에서 실행한다. \texttt{-\/-plan}은 경로만 확인하며 계산
성공을 뜻하지 않는다.

\begin{verbatim}
python tools/run_assignment.py A02 --plan
python tools/run_assignment.py A02 --run \
  --workspace learner --output ../A02_result.json
\end{verbatim}

\subsection{결과를
설명하기}\label{uxacb0uxacfcuxb97c-uxc124uxba85uxd558uxae30}

근은 함수의 전체 배율에 무관하다. 잔차와 좌표 오차를 같은 조건으로
취급했다면, 어느 종료 기준을 썼는지 먼저 찾아 비교하자.

\textbf{남길 증거.} bracket 폭과 반복 횟수 표; 성공 사례와 실패 상태 한
개; 함수 스케일 변화의 해석

\textbf{판단 기준.} 2점 예: F 스케일은 root를 바꾸지 않지만 residual
stopping은 달라짐을 실행으로 설명한다. 최초 예측의 정오보다, 결과를 본
뒤 이유를 고치고 다른 입력에도 적용하는지를 본다.

\subsection{막혔을 때}\label{uxb9c9uxd614uxc744-uxb54c}

\textbf{힌트 1.} 매 반복마다 유지되어야 하는 조건을 assert 대신 종이에
먼저 적는다.

\textbf{힌트 2.} 아주 큰 반대 부호 끝점에서 중점식의 overflow 가능성을
검사한다.

\subsection{다음 만남에 다시
묻기}\label{uxb2e4uxc74c-uxb9ccuxb0a8uxc5d0-uxb2e4uxc2dc-uxbb3buxae30}

다음 회차에는 연속인 삼차함수와 pole을 가진 함수에서 같은 부호 반전이
무엇을 뜻하는지 설명한다. 함수값을 보기 전에 필요한 전제를 말한다. 먼저
짧게 답한 뒤 본문과 자신의 기록을 열어 확인한다.

\textbf{범위.} pole 및 discontinuity에서는 부호변화만으로 근을 주장할 수
없다.

\end{document}
